\subsection{Financial Time-Series Extrapolation}
\label{sec:extrapolation}

\subsubsection{Problem Characteristics}
Financial Time-Series Extrapolation (FTSE) is formally defined as the task of generating future trajectories $\mathbf{x}_{t+1:T}$ given a historical context $\mathbf{x}_{1:t}$. Unlike traditional point forecasting, modern FTSE is approached as a probabilistic conditional generation problem: learning the conditional distribution $p(\mathbf{x}_{t+1:T} | \mathbf{x}_{1:t})$ to capture the inherent uncertainty and multi-modal nature of financial markets. This shift acknowledges that financial markets are stochastic systems with complex, time-varying dynamics. 

The problem exhibits several distinct characteristics that differentiate it from general sequence extrapolation:
\begin{enumerate}
    \item \textbf{Causality Constraint}: The generation must be strictly causal—future values depend solely on past observations without any look-ahead bias. This is analogous to autoregressive language modeling, but with continuous, real-valued tokens.
    \item \textbf{Non-Stationarity}: Financial time series exhibit time-varying statistical properties (e.g., changing volatility regimes, structural breaks). The conditional distribution $p(\mathbf{x}_{t+1:T} | \mathbf{x}_{1:t})$ itself evolves over time, requiring models to adapt to distributional shifts.
    \item \textbf{Low Signal-to-Noise Ratio}: Market data is dominated by noise, making it challenging to separate genuine predictive signals from random fluctuations. This often leads to overfitting in deterministic models.
    \item \textbf{High-Dimensional Dependencies}: In multi-variate settings, assets exhibit complex cross-sectional dependencies (e.g., lead-lag relationships, sector correlations). Modeling these dependencies is crucial for coherent joint generation.
\end{enumerate}

Formally, the objective is to learn a generator $G_\theta$ that maps the context $\mathbf{x}_{1:t}$ to a distribution over future sequences, often parameterized as a set of samples $\{\mathbf{x}_{t+1:T}^{(i)}\}_{i=1}^N$. The evaluation metrics extend beyond point-wise accuracy (e.g., MSE) to probabilistic scoring rules (e.g., CRPS, Energy Score) and domain-specific financial utility measures (e.g., portfolio returns, risk metrics).


\subsubsection{Classical and Hybrid Predictive Models}

Classical statistical frameworks have historically provided the mathematical bedrock for modeling financial time-series dynamics through rigorous probabilistic assumptions. A prominent example is the Auto-Regressive Integrated Moving Average (ARIMA) model \citep{box1970}, which explicitly formulates temporal dependencies by combining autoregressive (AR) processes, differencing (I) to achieve stationarity, and moving average (MA) components. For a given time series $y_t$, an $\text{ARIMA}(p, d, q)$ specification is mathematically expressed as:
\begin{equation}
\nabla^d y_t = c + \sum_{i=1}^p \phi_i \nabla^d y_{t-i} + \sum_{j=1}^q \theta_j \varepsilon_{t-j} + \varepsilon_t
\end{equation}
where $\nabla^d$ represents the $d$-th order difference operator, $\phi_i$ and $\theta_j$ are the AR and MA coefficients respectively, and $\varepsilon_t \sim \mathcal{N}(0, \sigma^2)$ acts as white noise. While ARIMA models are effective at isolating linear relationships, they inherently fail to encompass complex market phenomena such as volatility clustering and heavy-tailed distributions. To capture these time-varying variance dynamics, the Generalized Autoregressive Conditional Heteroskedasticity (GARCH) framework \citep{kim2021} was introduced. The conditional variance $\sigma_t^2$ in a $\text{GARCH}(p, q)$ model is iteratively defined by past squared innovations and past variances:
\begin{equation}
\sigma_t^2 = \omega + \sum_{i=1}^p \alpha_i \varepsilon_{t-i}^2 + \sum_{j=1}^q \beta_j \sigma_{t-j}^2
\end{equation}
where $\omega > 0$, $\alpha_i \ge 0$, and $\beta_j \ge 0$. Although extensions like EGARCH further accommodate asymmetric leverage effects resulting from negative shocks, these purely parametric methods struggle to scale to high-dimensional dependencies and intricate non-linear patterns.

To address the constraints of classical statistical models, the field transitioned toward deep learning architectures, particularly Recurrent Neural Networks (RNNs) and Long Short-Term Memory (LSTM) networks \citep{tsa1, qin2025, li2025ceemdan, caliskan2025, SGP-LSTM1}. LSTMs leverage a sophisticated gating mechanism—comprising forget ($f_t$), input ($i_t$), and output ($o_t$) gates alongside a cell state ($C_t$) and hidden state ($h_t$)—to effectively capture long-range temporal dependencies without suffering from severe vanishing gradients:
\begin{align}
f_t &= \sigma(\mathbf{W}_f \cdot[h_{t-1}, x_t] + b_f) \\
i_t &= \sigma(\mathbf{W}_i \cdot [h_{t-1}, x_t] + b_i) \\
\tilde{C}_t &= \tanh(\mathbf{W}_C \cdot[h_{t-1}, x_t] + b_C) \\
C_t &= f_t \odot C_{t-1} + i_t \odot \tilde{C}_t \\
o_t &= \sigma(\mathbf{W}_o \cdot [h_{t-1}, x_t] + b_o) \\
h_t &= o_t \odot \tanh(C_t)
\end{align}
However, when standard LSTMs are optimized using traditional pointwise loss functions such as the Mean Squared Error (MSE), defined as $\mathcal{L}_{MSE} = \frac{1}{n} \sum_{i=1}^n (y_i - \hat{y}_i)^2$, they often overfit the ubiquitous noise present in financial data.

Consequently, recent literature has pivoted towards sophisticated hybrid strategies that couple deep learning backbones with domain-specific signal processing or invariant learning techniques to enhance robustness. Denoising-based hybrids, for example, incorporate explicit signal decomposition to systematically separate underlying predictive trends from stochastic market noise. The PEC-W model \citep{caliskan2025} leverages the Discrete Wavelet Transform (DWT) to split raw input series into multi-resolution scale and wavelet representations:
\begin{equation}
x(t) = \sum_k c_{J,k} \phi_{J,k}(t) + \sum_{j=1}^J \sum_k d_{j,k} \psi_{j,k}(t)
\end{equation}
where $\phi$ and $\psi$ denote the scaling and wavelet functions. These denoised approximation coefficients provide a cleaner input signal for the LSTM, a concept similarly explored by CEEMDAN (Complete Ensemble Empirical Mode Decomposition with Adaptive Noise)-Informer-LSTM \citep{li2025ceemdan}, which decomposes the signal into intrinsic mode functions. Beyond denoising, feature engineering hybrids seek to automate the discovery of robust predictive inputs. The SGP-LSTM framework \citep{SGP-LSTM1} employs Symbolic Genetic Programming to iteratively evolve technical indicators, steering the search via a multi-objective fitness function that balances return, consistency, and predictive correlation (IC):
\begin{equation}
\text{Fitness} = \text{TopR} + \lambda_1 \times \text{Monotonicity} + \lambda_2 \times \text{IC}
\end{equation}
Finally, invariant learning hybrids target the problem of distributional shifts across varying market regimes. The Invariant Risk Minimization (IRM) approach proposed by \citet{cao2023irm} enforces the extraction of representations that hold true regardless of whether the market is in a bullish, bearish, or sideways phase, minimizing an objective that penalizes environmental variance:
\begin{equation}
\mathcal{L}_{IRM} = \underbrace{R^e(\mathbf{w} \cdot \Phi)}_{\text{Empirical risk}} + \gamma \underbrace{\left\| \nabla_{\mathbf{w}|_{\mathbf{w}=1.0}} R^e(\mathbf{w} \cdot \Phi) \right\|_2^2}_{\text{Invariant penalty}}
\end{equation}
where $\Phi$ acts as the representation function and $\mathbf{w}$ is a linear classifier. Despite these multifaceted improvements, hybrid predictive architectures ultimately generate deterministic point forecasts, fundamentally limiting their capacity to model the full predictive distribution essential for robust financial scenario generation.

\subsubsection{Generative Models for Extrapolation}

Recognizing the intrinsic limitations of deterministic forecasting, the paradigm of generative models shifts the focus toward learning the underlying data distribution $p(\mathbf{x})$ or the conditional probability $p(\mathbf{x}_{t+1:T}|\mathbf{x}_{1:t})$. This transition allows practitioners to sample multiple plausible future trajectories, accommodating the multi-modal uncertainty inherent in financial markets. Early efforts in this domain heavily utilized Generative Adversarial Networks (GANs) \citep{Time-GAN, naritomi2020data}, which deploy a generator $G$ and a discriminator $D$ locked in a competitive minimax game. The generator synthesizes sequences from historical context $\mathbf{x}_{1:t}$ and a latent noise vector $\mathbf{z}$, while the discriminator attempts to distinguish real historical samples from the generated distributions:
\begin{equation}
\min_G \max_D \mathbb{E}_{\mathbf{x} \sim p_{data}}[\log D(\mathbf{x})] + \mathbb{E}_{\mathbf{z} \sim p_z}[\log(1 - D(G(\mathbf{z})))]
\end{equation}
While conceptually elegant, GANs are notoriously difficult to train on volatile financial data, frequently suffering from mode collapse and unstable convergence. To mitigate these issues, specialized adaptations have been devised. WaveGAN \citep{dai2024wavegan} incorporates a multi-resolution DWT denoising mechanism combined with a Wasserstein GAN equipped with Gradient Penalty (WGAN-GP) \citep{wang2024denoiser}. The WGAN-GP loss enforces a Lipschitz constraint to prevent vanishing gradients:
\begin{equation}
\mathcal{L}_{\text{WGAN-GP}} = \mathbb{E}_{\tilde{\mathbf{x}} \sim p_g}[D(\tilde{\mathbf{x}})] - \mathbb{E}_{\mathbf{x} \sim p_{data}}[D(\mathbf{x})] + \lambda \mathbb{E}_{\hat{\mathbf{x}} \sim p_{\hat{\mathbf{x}}}}[(\|\nabla_{\hat{\mathbf{x}}} D(\hat{\mathbf{x}})\|_2 - 1)^2]
\end{equation}
where $\hat{\mathbf{x}}$ represents a convex interpolation of real and generated instances. Concurrently, models like N-BEATS-GAN \citep{dai2026nbeats} fuse interpretable basis expansions \citep{oreshkin2020nbeats} with an adversarial framework, modeling future values as a sum of trend and seasonality bases $\hat{\mathbf{y}} = \sum_{i=1}^N \text{Basis}_i(\boldsymbol{\theta}_i)$. While these modifications enhance stability and interpretability, GANs often yield over-smoothed samples that fail to encapsulate the sharp, multi-modal nature of empirical return distributions.

To achieve higher fidelity in distribution matching, Denoising Diffusion Probabilistic Models (DDPMs) \citep{ho2020ddpm} have emerged as a more stable and expressive alternative. Diffusion frameworks systematically destroy structured data by injecting Gaussian noise over $T$ forward steps—$q(\mathbf{x}_t|\mathbf{x}_{t-1}) = \mathcal{N}(\mathbf{x}_t; \sqrt{1 - \beta_t}\mathbf{x}_{t-1}, \beta_t \mathbf{I})$—and then learn a parameterized reverse process $p_\theta(\mathbf{x}_{t-1}|\mathbf{x}_t) = \mathcal{N}(\mathbf{x}_{t-1}; \boldsymbol{\mu}_\theta(\mathbf{x}_t, t), \sigma_t^2 \mathbf{I})$ to reconstruct the signal. The optimization simplifies to predicting the incrementally added noise $\boldsymbol{\epsilon}$:
\begin{equation}
\mathcal{L}_{simple} = \mathbb{E}_{t, \mathbf{x}_0, \boldsymbol{\epsilon}}[\|\boldsymbol{\epsilon} - \boldsymbol{\epsilon}_\theta(\mathbf{x}_t, t)\|^2]
\end{equation}
For context-conditioned financial generation, TimeGrad \citep{rasul2021timegrad} integrates an RNN-based encoder with a score-based diffusion decoder, generating future states via Langevin dynamics where the context encoding $\mathbf{c}$ steers the reverse sampling:
\begin{equation}
\mathbf{x}_{t-1} = \frac{1}{\sqrt{\alpha_t}} \left( \mathbf{x}_t - \frac{\beta_t}{\sqrt{1 - \bar{\alpha}_t}} \boldsymbol{\epsilon}_\theta(\mathbf{x}_t, t, \mathbf{c}) \right) + \sigma_t \mathbf{z}
\end{equation}
Addressing the limitations of recurrent architectures in capturing extended dependencies, TimeDiT \citep{cao2024timedit} replaces the RNN with a Diffusion Transformer. Utilizing patch embeddings and multi-head attention, TimeDiT operates on a $\mu$-parameterized loss function, explicitly aligning the predicted mean of the reverse process with its theoretical target:
\begin{equation}
\mathcal{L} = \sum_{t=1}^T \mathbb{E}_{q(\mathbf{x}_t^{tar}|\mathbf{x}_0^{con})} \left[ \|\boldsymbol{\mu}(\mathbf{x}_t^{tar}, \mathbf{x}_0^{con}) - \boldsymbol{\mu}_\theta(\mathbf{x}_t^{tar}, t|\mathbf{x}_0^{con})\|^2 \right]
\end{equation}
An alternative diffusion formulation leverages continuous-time score matching, optimizing a score network to approximate the gradient of the log-density:
\begin{equation}
\mathcal{L}_{score} = \mathbb{E}_{t, \mathbf{x}_0, \boldsymbol{\epsilon}} \left[ \left\| \boldsymbol{\epsilon} - \sigma_t \nabla_{\mathbf{x}_t} \log p_\theta(\mathbf{x}_t) \right\|^2 \right]
\end{equation}
While diffusion and score-based models generate high-fidelity distributions with sharp multi-modal properties, their iterative sampling procedures introduce significant inference latency, presenting an ongoing challenge for high-frequency financial applications.



\subsubsection{Time Series Foundation Models}

Inspired by the transformative success of Large Language Models (LLMs) in natural language processing \citep{Survey1, Survey2, tsa2}, current research is heavily invested in adapting the "pre-train, then fine-tune" paradigm to financial time-series extrapolation. These Foundation Models are trained on massive, heterogeneous datasets to derive generalized temporal representations, effectively enabling zero-shot inference and cross-domain adaptation. The primary challenge in this domain is bridging the continuous numerical space of financial data with the discrete, tokenized vocabulary natively processed by transformer-based LLMs. Tokenization-based fine-tuning strategies directly address this by casting forecasting as a next-token prediction task. For instance, Chronos \citep{ansari2024chronos} scales time-series observations and maps them into a fixed dictionary of tokens using quantile binning:
\begin{equation}
\text{token}(\mathbf{x}) = \arg\min_i |x - q_i|
\end{equation}
where $q_i$ represent pre-computed quantiles. While this approach benefits from robust zero-shot generalization, the inherent discretization process sacrifices fine-grained numerical precision. Conversely, patch-based fine-tuning techniques, as exemplified by OneFitsAll \citep{zhou2023onefitsall}, bypass discretization by segmenting the time-series into sequences of sub-series patches. Each patch is projected into the high-dimensional embedding space of the LLM via a linear transformation:
\begin{equation}
\mathbf{E}_{patch} = \mathbf{W}_p \mathbf{x}_{patch} + \mathbf{b}_p
\end{equation}
These continuous patch embeddings are then prepended as visual-like tokens, fully preserving numerical fidelity while seamlessly leveraging the LLM's cross-patch attention mechanisms.

As illustrated in Figure \ref{fig:llm_comparison}, beyond conventional fine-tuning, sophisticated model editing and reprogramming approaches have emerged to align disparate modalities using lightweight adapters, averting the computational burden of full-parameter updates. StockTime \citep{wang2024stocktime} enriches time-series representations by encoding high-frequency tick data alongside handcrafted textual summaries (e.g., "bullish engulfing pattern") via a dual-encoder. Similarly, frameworks like LLM4FTS \citep{liu2025llm4fts} and LLM-PS \citep{tang2025llmps} utilize Dynamic Time Warping (DTW) and text-translation modules to bridge the semantic divide between raw price trajectories and natural language descriptions. A more fundamental alignment is achieved through Prompt-as-Prefix reprogramming, popularized by Time-LLM \citep{jin2024timellm}. This technique steers frozen LLMs by projecting time-series patches onto a set of learned "text prototypes" $\mathbf{P} \in \mathbb{R}^{M \times d}$ using cross-attention dynamics:
\begin{align}
\alpha_{ij} &= \text{softmax} \left(\frac{(\mathbf{W}_q \mathbf{x}_i) \cdot (\mathbf{W}_k \mathbf{p}_j)^\top}{\sqrt{d}}\right) \\
\tilde{\mathbf{x}}_i &= \sum_j \alpha_{ij} \mathbf{p}_j
\end{align}
The resulting reprogrammed patches act as informative prefixes for the frozen language model. Expanding upon this, the CALF architecture \citep{liu2025calf} employs a dual-branch structure fortified by consistency losses to strictly minimize the discrepancy between textual and temporal feature distributions.

In parallel, researchers are developing purpose-built forecasting architectures that integrate domain-specific inductive biases directly into the foundation model backbone. TimeHF \citep{qi2025timehf} introduces a time-series policy optimization (TSPO) framework that aligns generative outputs with specific financial human preferences, such as maximizing risk-adjusted returns, using reinforcement learning from human feedback (RLHF):
\begin{equation}
\mathcal{L}_{TSPO} = -\mathbb{E}_{\tau \sim \pi_\theta}[R(\tau)] + \beta \text{KL}[\pi_\theta \| \pi_{ref}]
\end{equation}
where $\pi_\theta$ denotes the policy (generative model) parameterized by $\theta$, $\tau \sim \pi_\theta$ is a sampled generation trajectory, $R(\tau)$ is the reward function that scores trajectory quality according to human financial preferences, $\beta$ controls the strength of the KL-divergence regularization, and $\pi_{\text{ref}}$ is the frozen reference policy that prevents catastrophic drift from pre-trained behavior.
FinCast \citep{zhu2025fincast} attacks the multi-scale nature of market data using a sparse Mixture-of-Experts (MoE) backbone enhanced with learnable frequency embeddings. To ensure the model delivers both accurate point estimates and reliable probabilistic bounds, it optimizes a joint Point-Quantile loss:
\begin{align}
\mathcal{L}_{PQ} &= \lambda_1 \mathcal{L}_{point} + \lambda_2 \mathcal{L}_{quantile} \\
\mathcal{L}_{quantile} &= \sum_q \rho_q(y - \hat{y}_q), \quad \rho_q(u) = u(q - \mathbb{I}(u < 0))
\end{align}
where $\mathcal{L}_{point}$ handles deterministic accuracy and the pinball loss $\mathcal{L}_{quantile}$ governs risk coverage, yielding highly robust zero-shot performance across diverse asset classes.

\subsubsection{Emerging Trends}
Moving forward, the field is poised to tackle several critical bottlenecks that limit the operationalization of generative models in live financial systems. A primary objective is the development of unified hybrid architectures that can seamlessly amalgamate the low-latency inference speeds of discriminative models with the distributional robustness of diffusion frameworks, potentially leveraging sophisticated knowledge distillation techniques. Concurrently, there is an urgent need to redefine evaluation protocols; existing metrics such as the Continuous Ranked Probability Score (CRPS) must be augmented with practical financial utility measures, including risk-adjusted returns and maximum drawdowns, to ensure models align with authentic trading objectives. Furthermore, future iterations of generative models must incorporate causal constraints that respect intricate market microstructures, fundamentally eliminating the look-ahead biases that persistently plague synthetic market simulators. Finally, advancing computationally efficient fine-tuning protocols will be paramount to ensure that massive foundation models can rapidly adapt to sudden structural regime shifts without succumbing to catastrophic forgetting.

\subsubsection{Summarization}
The progressive evolution of financial time-series extrapolation models demonstrates a clear overarching trajectory in machine learning: transitioning from deterministic, discriminative point predictors to expressive generative distribution learners, and ultimately culminating in large-scale foundation models. As synthesized in Table \ref{tab:extrapolation_comparison}, each methodological paradigm presents distinct trade-offs between inferential speed, distributional fidelity, and generalization capability. Classical and hybrid predictive models, such as ARIMA and LSTM-based architectures, retain their dominance in latency-sensitive environments like high-frequency execution algorithms, where rapid point forecasting is paramount. However, their structural inability to quantify tail-risk and multimodal uncertainty severely restricts their utility in comprehensive risk management. To resolve this, generative architectures—notably GANs and diffusion models—excel in scenario generation and stress testing by synthesizing highly realistic, multi-modal market trajectories. Recent studies often report that diffusion models yield better calibrated distributions than GAN baselines, although this gain typically comes at the cost of substantially higher computational burden during iterative reverse-sampling. Time-series foundation models have recently become an important direction by exhibiting unprecedented zero-shot forecasting capabilities, enabling the deployment of predictive systems across previously unseen assets or shifting market regimes. Despite these breakthroughs, foundation models still grapple with persistent challenges regarding temporal alignment, cross-modal semantic gaps, and the propensity to hallucinate implausible financial patterns over extended horizons.

